\documentclass{letter}
\usepackage{dirtytalk}
\usepackage[letterpaper, margin=2.5cm]{geometry}

\signature{Dr. Jill Emmerzaal}
\address{McGill University \\ 475 Pine Avenue West \\ Montreal \\ H2W 1S4, Quebec \\ Canada}
\begin{document}
\longindentation=0pt

\begin{letter}{}
\opening{Dear Editors-in-Chief,}

Hereby, we submit the manuscript \say{Gait dynamics are affected by commonly occurring outdoor surfaces and sex in healthy adults} for publication as an original article in \textit{Journal of Biomechanics}. All authors confirm that they have made substantial contributions to the conception and design of the study, analysis and interpretation of data, drafting the article and revising it critically, and give final approval of the version submitted here. An abstract of this work has been submitted for the ISB2025 conference to be held in Stockholm later this year. 

To gain insights into the impact of surfaces on outdoor walking for males and females this study aimed to investigate the the effect of surface and sex on non-linear walking dynamics on commonly occurring outdoor surfaces.
We used an existing, publicly available dataset of thirty healthy young adults equipped with five inertial measurement units (IMUs) on the lower limbs and lower back. Data were extracted from the lower back, left, and right shank for seven surfaces: flat-even, banked (right, left), irregular (cobblestone, grass), and sloped (down, up), with six trials per surface. The IMU signals were concatenated into a single time series, and from the data of the lower trunk IMU we calculated non-linear gait dynamics measures: sample entropy, log dimensionless jerk (LDLJ), short and long term maximum Lyapunov exponent (LyE$_{s}$ and LyE$_{l}$, respectively), and step and stride symmetry.

We found that walking surfaces impacted many aspects of gait dynamics (i.e., predictability, smoothness, symmetry, and stability), independent of sex. Broadly, when compared to flat walking, reduced LDLJ (smoothness) and increased long-term Lyapunov exponent (complexity) were the most consistent findings, while sloped-down, cobble-stones and banked right surfaces appear to challenge gait dynamics the most. No significant sex differences were found in gait dynamics; however, males exhibited a non-significant trend towards lower short-term Lyapunov exponent, suggesting less step-to-step variability. Significant interactions between sex and surface were observed for Sample Entropy and stride symmetry on banked right surfaces, with females showing decreased sample entropy and increased stride symmetry compared to males.

We are convinced that this manuscript is suitable for publication in \textit{Journal of Biomechanics} as the results bring us one step closer to understanding the demands of outdoor surfaces on walking dynamics. To the best of our knowledge, we are unique on two fronts, first, by incorporating a large number of outdoor surfaces we gain comprehensive insights into the demands of various outdoor surfaces on gait dynamics. Second, by incorporating sex in our analysis we showed that greater consideration of sex differences in gait studies is necessary to enhance understanding of movement adaptation strategies.

Thank you for your time and consideration and I remain available should you require further information. I look forward to your reply.

\closing{Yours sincerely,}

\end{letter}
\end{document}

